% Démonstration de la formule locale de la divergence discrète.
% Ce fichier peut être inclus à la place de l'environnement `demo` vide.

Soit $u^\ensddfvplat \in \R^\ensddfvplat$. Par définition de la
divergence discrète comme opposé de l'adjoint du gradient,
\begin{equation}\label{eq:preuve_div_adjonction}
    \prodscalddfv{u^\ensddfvplat}{\div^\ensddfvplat\xi^\ensdiamantplat}
    =-
    \prodscalD{\grad^\ensdiamantplat u^\ensddfvplat}{\xi^\ensdiamantplat}.
\end{equation}
Développons le membre de droite diamant par diamant. Pour
$\diamantplat{D}_{\sigma,\dual{\sigma}}$, la formule du gradient donne
\[
\grad^{\diamantplat{D}}u^\ensddfvplat
=\frac{1}{\sin\angle{\diamantplat{D}}}
\left(
 \frac{u_{\primalplat{L}}-u_{\primalplat{K}}}{\mesure(\dual{\sigma})}
 \vect{\nu}_{\sigma\primalplat{K}}
 +
 \frac{u_{\mailledualeplate{L}}-u_{\mailledualeplate{K}}}{\mesure(\sigma)}
 \vect{\nu}_{\dual{\sigma}\mailledualeplate{K}}
\right).
\]
Or
\[
\mesure(\diamantplat{D})
=\frac12\mesure(\sigma)\mesure(\dual{\sigma})
 \sin\angle{\diamantplat{D}}.
\]
Par conséquent,
\begin{align*}
-\prodscalD{\grad^\ensdiamantplat u^\ensddfvplat}{\xi^\ensdiamantplat}
={}&-\frac12\sum_{\diamantplat{D}_{\sigma,\dual{\sigma}}\in\ensdiamantplat}
 \mesure(\sigma)
 (u_{\primalplat{L}}-u_{\primalplat{K}})
 \xi^{\diamantplat{D}}\cdot\vect{\nu}_{\sigma\primalplat{K}}
\\
&-\frac12\sum_{\diamantplat{D}_{\sigma,\dual{\sigma}}\in\ensdiamantplat}
 \mesure(\dual{\sigma})
 (u_{\mailledualeplate{L}}-u_{\mailledualeplate{K}})
 \xi^{\diamantplat{D}}\cdot
 \vect{\nu}_{\dual{\sigma}\mailledualeplate{K}}.
\end{align*}

Regroupons maintenant les termes qui multiplient une même valeur
$u_{\primalplat{K}}$. Sur une arête $\sigma$ de $\primalplat{K}$, la
normale $\vect{\nu}_{\sigma\primalplat{K}}$ est la normale extérieure à
$\primalplat{K}$. Si $\primalplat{K}$ est désignée comme la seconde
maille du diamant, l'inversion simultanée de la différence des valeurs
et de la normale donne le même terme. Le coefficient de
$u_{\primalplat{K}}$ dans le membre de droite de
\eqref{eq:preuve_div_adjonction} est donc
\[
\frac12
\sum_{\diamantplat{D}_{\sigma,\dual{\sigma}}
      \in\ensdiamantplat_{\primalplat{K}}}
\mesure(\sigma)\,
\xi^{\diamantplat{D}}\cdot\vect{\nu}_{\sigma\primalplat{K}}.
\]
De même, le coefficient de $u_{\mailledualeplate{K}}$ vaut
\[
\frac12
\sum_{\diamantplat{D}_{\sigma,\dual{\sigma}}
      \in\ensdiamantplat_{\mailledualeplate{K}}}
\mesure(\dual{\sigma})\,
\xi^{\diamantplat{D}}\cdot
\vect{\nu}_{\dual{\sigma}\mailledualeplate{K}}.
\]

D'autre part, puisque
$\poids{\primalplat{K}}=\mesure(\primalplat{K})/2$ et
$\poids{\mailledualeplate{K}}=\mesure(\mailledualeplate{K})/2$, le
membre de gauche de \eqref{eq:preuve_div_adjonction} est
\begin{align*}
&\frac12\sum_{\primalplat{K}\in\ensprimalplat}
 \mesure(\primalplat{K})u_{\primalplat{K}}
 \div^{\primalplat{K}}\xi^\ensdiamantplat
\\
&\qquad+
\frac12\sum_{\mailledualeplate{K}\in\ensdualplat}
 \mesure(\mailledualeplate{K})u_{\mailledualeplate{K}}
 \div^{\mailledualeplate{K}}\xi^\ensdiamantplat.
\end{align*}
L'identité étant vraie pour tout $u^\ensddfvplat$, l'identification des
coefficients de ses composantes fournit
\[
\div^{\primalplat{K}}\xi^\ensdiamantplat
=\frac{1}{\mesure(\primalplat{K})}
 \sum_{\diamantplat{D}_{\sigma,\dual{\sigma}}
       \in\ensdiamantplat_{\primalplat{K}}}
 \mesure(\sigma)\,
 \xi^{\diamantplat{D}}\cdot\vect{\nu}_{\sigma\primalplat{K}},
\]
et
\[
\div^{\mailledualeplate{K}}\xi^\ensdiamantplat
=\frac{1}{\mesure(\mailledualeplate{K})}
 \sum_{\diamantplat{D}_{\sigma,\dual{\sigma}}
       \in\ensdiamantplat_{\mailledualeplate{K}}}
 \mesure(\dual{\sigma})\,
 \xi^{\diamantplat{D}}\cdot
 \vect{\nu}_{\dual{\sigma}\mailledualeplate{K}},
\]
ce qui est précisément le résultat annoncé.
